\section{Cross-Cutting Gaps}\label{sec:synthesis}

Read side by side, the open problems of the six functions repeat themselves in different vocabularies. Take
a viewport predictor never tested in a closed loop, an offloading controller that never left the
simulator, and an allocation scheme whose QoE objective was never shown to a user. They are three
instances of one shortfall. The atlas verified \nxrSynGapsTotal{} gaps across its \nxrDeepCells{} cells,
\nxrSynGapsHigh{} of them with high confidence. Nine themes account for \nxrSynCovered{} of them
(\nxrSynCoveredShare{}), including \nxrSynCoveredHigh{} of the high-confidence ones
(\cref{tab:themes}). A theme was kept only if it held gaps from at least six cells. We group the nine
themes into three families.

\begin{table}[t]
\centering\small
\caption{Nine cross-cutting themes. Gaps: verified cell gaps that instance the theme (high-confidence in
brackets). Cells: atlas cells in which the theme recurs. A gap can belong to more than one
theme.}\label{tab:themes}
\setlength{\tabcolsep}{4pt}
\begin{tabularx}{\linewidth}{@{}l>{\raggedright\arraybackslash}p{4.6cm}rr>{\raggedright\arraybackslash}X@{}}
\toprule
& Theme & Gaps & Cells & Indicative evidence\\
\midrule
\multicolumn{5}{@{}l}{\emph{Evidence foundations}}\\
T1 & No open traces, datasets or benchmarks & \nxrSynDataGapGaps{} (\nxrSynDataGapHigh{}) & \nxrSynDataGapCells{} & \nxrNDatasets{} dataset papers; \nxrSynReleaseShare{} of abstracts mention released code or data\\
T2 & Evaluation stops at the simulator & \nxrSynSimOnlyGaps{} (\nxrSynSimOnlyHigh{}) & \nxrSynSimOnlyCells{} & \nxrSynLearnedReal{} of \nxrSynLearnedN{} learning papers tested on a real network\\
T3 & Standards are not the evaluation substrate & \nxrSynStandardsGaps{} (\nxrSynStandardsHigh{}) & \nxrSynStandardsCells{} & \nxrSynThreeGPPXRN{} XR abstracts name 3GPP; none evaluated with users\\
\addlinespace
\multicolumn{5}{@{}l}{\emph{Objectives}}\\
T4 & Network metrics not linked to experience & \nxrSynNoQoEGaps{} (\nxrSynNoQoEHigh{}) & \nxrSynNoQoECells{} & \nxrSynQoEtermN{} abstracts invoke QoE; \nxrSynQoEtermUsers{} test it with users\\
T5 & Requirement figures inherited, not derived & \nxrSynRequirementsGaps{} (\nxrSynRequirementsHigh{}) & \nxrSynRequirementsCells{} & 1~ms target in \nxrSynOneMsN{} abstracts, \nxrSynOneMsUsersN{} tested with users\\
\addlinespace
\multicolumn{5}{@{}l}{\emph{Blind spots}}\\
T6 & Uplink and sending side neglected & \nxrSynUplinkGaps{} (\nxrSynUplinkHigh{}) & \nxrSynUplinkCells{} & ``uplink'' in \nxrSynUplinktermShare{} of abstracts, \nxrSynUplinktermSimAna{} of them simulated\\
T7 & Multi-user scale untested & \nxrSynScaleGaps{} (\nxrSynScaleHigh{}) & \nxrSynScaleCells{} & largest shared-world test: 225 clients \cite{W4400277138}\\
T8 & Security and privacy not costed & \nxrSynSecurityGaps{} (\nxrSynSecurityHigh{}) & \nxrSynSecurityCells{} & \nxrSynSecLatUsersN{} of \nxrSynSecLatN{} security--latency papers tested with users\\
T9 & Transport under real impairments & \nxrSynTransportGaps{} (\nxrSynTransportHigh{}) & \nxrSynTransportCells{} & \nxrSynXLFourSN{} corpus papers apply L4S\\
\bottomrule
\end{tabularx}
\end{table}

\subsection{Evidence foundations}
The first family concerns what claims are tested against. \emph{Data are missing} (T1). There are no
public traffic traces for AR offloading, avatar streams, twin-state updates, holographic calls or haptic
flows. No 360-degree dataset records head motion, network conditions and ratings in the same session,
and no latency benchmark couples motion, network delay and photodiode-measured MTP. Even simple
baselines are often missing: a persistence predictor reached a hit rate of 0.8355 against 0.8402 for a
sequence-to-sequence model on the same data \cite{W2963428739}. \emph{Evaluation stops at the simulator}
(T2). Whole technology lines, including reconfigurable surfaces carrying XR, terahertz XR links and
metaverse resource markets, have never met hardware. Where learned controllers were tested on unseen
users or content, they degraded \cite{W4403826584,W4400409717}. \emph{Standards are not used as the
substrate} (T3). 3GPP XR traffic models appear almost only in radio simulations. When one was checked
against measurement, it underestimated cloud-VR bitrate by 25\% \cite{W4399340506}. OMAF, MPEG Immersive
Video, IEEE~1918.1.1 and ITU-T P.1203/P.1204 are rarely the baseline configuration. The three gaps share
a cause. Collecting traces and building testbeds is costly, and venues that publish algorithms rarely
credit either.

\subsection{Objectives}
The second family concerns what is optimised. \emph{Network metrics are not linked to experience} (T4).
This is the largest theme. Most papers that optimise ``QoE'' optimise a formula whose weights were
chosen rather than fitted. Where users were asked, the answers diverged from the proxies. Packet loss
outweighed delay in VR cloud gaming \cite{W4404916373}. Users of collaborative VR could not tell forms
of latency apart \cite{W4392600587}. CloudVR users did not notice 20~ms of added latency
\cite{W2896076858}, while offloading papers optimise to budgets of 5--13~ms. \emph{Requirements are
inherited} (T5). Vision papers state 100--250~Mbps per metaverse user \cite{W4313291084,W4312940001},
while a measured commercial social-VR platform used 2--3~Mbps downstream \cite{W4294643390}. Holographic
requirements range from 5--100~Mbps for compressed calls \cite{W4312662428} to 6.63~TBps
\cite{W4300960574}. Haptic users tolerated up to 200~ms with adaptive control \cite{W4200154132}, against
the 1~ms figure that motivates much tactile-internet work. A requirement that names neither the
representation, the codec, the loop nor the percentile cannot be checked.

\subsection{Blind spots}
The third family concerns what is left out. XR sends a continuous upstream, yet the \emph{uplink} (T6)
is studied far less than the downlink, and the measured cases are pessimistic. On public LTE, multi-user
AR had a median end-to-end latency of 12.5~s, with the radio access network contributing 31.2\%
\cite{W3047439371}. With three headsets on Wi-Fi~6, uplink OFDMA delayed tracking data
\cite{W4396853623}. \emph{Scale} (T7) is untested. Measured social platforms scale linearly because they
forward all avatar data \cite{W4307020344}, and controlled QoE studies use pairs of users. \emph{Security
is not costed} (T8). Attacks are measured and effective: keystrokes were inferred from avatar motion with
97.62\% accuracy \cite{W4398796475}. Defences, however, are seldom priced in frame latency. \emph{Transport}
(T9) has not followed the move from buffered streaming to interactive delivery. Users rank loss above
delay \cite{W3095344018,W4404916373}, yet forward error correction and L4S are rare. An L4S pipeline for
volumetric conferencing converged 45\% faster than Google Congestion Control \cite{W4408870658}, but in a
controlled bottleneck.

\begin{figure}[t]
\centering
\begin{tikzpicture}[font=\sffamily\scriptsize,
  th/.style={draw=black!55,rounded corners=2pt,align=center,text width=2.35cm,minimum height=8mm,inner sep=2pt,fill=white},
  root/.style={th,fill=gapcolor!8,draw=gapcolor!70!black,line width=0.7pt},
  arr/.style={-{Stealth[length=1.6mm]},black!55,line width=0.55pt}]
  \node[root] (t1) at (0,0.55) {\textbf{T1} Open traces\\and datasets};
  \node[root] (t4) at (0,-0.75) {\textbf{T4} QoE linked to\\network metrics};
  \node[th] (t2) at (3.6,1.25) {\textbf{T2} Sim-to-real\\validation};
  \node[th] (t3) at (3.6,0.2) {\textbf{T3} Standards as\\substrate};
  \node[th] (t5) at (3.6,-0.85) {\textbf{T5} Derived\\requirements};
  \node[th] (t6) at (7.2,1.25) {\textbf{T6} Uplink};
  \node[th] (t7) at (7.2,0.2) {\textbf{T7} Multi-user scale};
  \node[th] (t8) at (7.2,-0.85) {\textbf{T8} Costed security};
  \node[th] (t9) at (10.8,0.2) {\textbf{T9} Transport under\\real impairments};
  \foreach \a/\b in {t1/t2,t1/t3,t1/t5,t4/t5,t4/t3}{\draw[arr] (\a.east) -- (\b.west);}
  \draw[arr] (t1.north) |- (t6.west);
  \draw[arr] (t2.east) -- (t7.west);
  \draw[arr] (t4.south) |- ++(0,-0.55) -| (t8.south);
  \draw[arr] (t6.east) -- (t9.north west);
  \draw[arr] (t7.east) -- (t9.west);
  \draw[arr] (t8.east) -- (t9.south west);
  \node[font=\sffamily\tiny,text=black!60,align=left,anchor=west] at (9.6,-0.95) {arrow: the target theme\\needs the source's artefact};
\end{tikzpicture}
\caption{Dependencies among the nine themes. Open traces (T1) and a validated mapping from network
metrics to experience (T4) are the roots. Without them, simulations cannot be validated, requirements
cannot be derived, and allocation, transport and security results cannot be compared in units that matter
to users.}\label{fig:themes}
\end{figure}

\subsection{How the themes depend on each other}
The themes are not independent (\cref{fig:themes}). Two are roots. Without open traces (T1),
simulations cannot be validated (T2), requirements cannot be derived (T5), standard traffic models cannot
be checked (T3), and uplink and multi-user workloads cannot even be described (T6, T7). Without a
validated mapping from network metrics to experience (T4), results on allocation, transport and security
(T8, T9) cannot be compared in units that matter to users. The two smaller clusters outside the themes
point the same way. Device-energy savings from offloading rest on models, and the few measurements
disagree: one standalone headset drained its battery 15\% faster when offloading wirelessly
\cite{W4362489903}, while an edge VR testbed found lower client power \cite{W4296339041}. Neural
scene representations have no tested rate adaptation, cache or split point. Both clusters need the same
two roots. This ordering shapes the agenda that follows.
